\documentclass{article}
\usepackage[utf8]{inputenc}
\usepackage[ukrainian,english]{babel}
\usepackage{amsmath,amssymb,amsfonts}
\usepackage{graphicx}
\usepackage{hyperref}
\usepackage{listings}
\usepackage{booktabs}
\usepackage{xcolor}

\lstset{
    keywordstyle=\color{blue},
    stringstyle=\color{red},
    basicstyle=\ttfamily\scriptsize,
    breakatwhitespace=false,
    keepspaces=true,
    captionpos=b,
    showspaces=false,
    showstringspaces=false,
    showtabs=false,
    tabsize=2
}

\title{\textbf{Audit Accountable Infrastructure:}\\Formal Taxonomy, Multi-Plane Cryptographic Logs, and NIST SP 800-53 Attestation}
\author{\textbf{Synrc Research Center} \\ \texttt{https://synrc.com}}
\date{\today}

\begin{document}

\maketitle

\begin{abstract}
This document presents the formal specification, architecture, and theoretical foundation of the \texttt{synrc/au}
library for high-assurance operating environments based on the ERP.1 three-plane architecture. Built upon standard
Erlang/OTP primitives and standard NIST cryptography (ECDSA P-384, SHA-384, HMAC-SHA-512), \texttt{synrc/au}
enforces strict privilege segregation, fail-secure append-only event logging, deterministic CRDT Merkle-log
merging across multi-DC deployments, cold-storage log archiving, and complete compliance with
NIST SP 800-53 Rev. 5 Audit and Accountability (AU) controls.
\end{abstract}

\newpage
\tableofcontents

\newpage
\section{Introduction \& Motivation}

High-integrity critical systems (defense appliances, state registries, electronic healthcare,
infrastructure control systems) demand rigorous mathematical assurance, non-repudiation, and
architectural privilege isolation. Traditional POSIX models featuring monolithic \texttt{root}
processes violate Zero Trust Architecture principles by exposing broad ambient authority.

The ERP.1 architecture reduces the Trusted Computing Base (TCB) to a minimal cryptographic core
running on Erlang/OTP on UKI Alpine Linux. The \texttt{synrc/au} library isolates audit generation,
aggregation, and signature capabilities across three execution planes, guaranteeing that audit
logs remain immutable, cryptographically verifiable, and fail-secure under all operating conditions.

\section{Audit and Accountability}

\subsection{ERP.1 Three-Plane Isolation Model}

Following the fundamental OS.1 taxonomy defined in \texttt{synrc/os/os.txt},
the execution environment is partitioned into three distinct planes:

\begin{figure}[ht!]
\centering
\begin{lstlisting}
+-----------------------------------------------------------------------------+
| Security Admin BEAM   [S]                   (Trusted Security Services)     |
|  * Highest privilege, minimal TCB (ca, tss, sealed kvs)                     |
|  * Owns long-term ECDSA P-384 keys, HSM/TPM, root & intermediate CAs        |
|  * Never opens network sockets or mounts filesystem volumes                 |
|  * Signs critical audit records & periodic Merkle checkpoints               |
+-----------------------------------------------------------------------------+
                : sealed, capability-limited length-prefixed channel
+-----------------------------------------------------------------------------+
| System Admin BEAM     [C]                   (Privileged Control Plane)      |
|  * Sole owner of kernel operations, seats, network & capability broker      |
|  * Maintains local append-only hash-chain (SHA-384 + HMAC-SHA-512)          |
|  * Synchronous fail-secure commits (AU-5)                                   |
|  * Forwards logs to SIEM (Wazuh)                                            |
+-----------------------------------------------------------------------------+
                : capability-granted, authenticated bridge
+-----------------------------------------------------------------------------+
| Application BEAM(s)   [A]                   (Untrusted Domain & Sessions)   |
|  * Host user interfaces, session logic, and media pipelines                 |
|  * No access to long-term keys or ambient storage                           |
|  * Generates local events forwarded exclusively via capability bridge       |
+-----------------------------------------------------------------------------+
\end{lstlisting}
\caption{ERP.1 Three-Plane Architecture Isolation Diagram.}
\end{figure}

\subsubsection{Plane Descriptions}
\begin{enumerate}
    \item \textbf{Security Admin BEAM [S]}: Houses \texttt{synrc/ca}, \texttt{synrc/tss}, and sealed \texttt{synrc/kvs}. Long-term private keys never leave this plane. It provides RFC 3161 timestamps and ECDSA signatures for critical events.
    \item \textbf{System Admin BEAM [C]}: Runs \texttt{audit\_core}. Handles capability broker passing via \texttt{SCM\_RIGHTS}, maintains the append-only record chain, applies HMAC-SHA-512 integrity tags, and publishes periodic checkpoints.
    \item \textbf{Application BEAM [A]}: Runs \texttt{audit\_client}. Untrusted workloads operating under strict Landlock, seccomp-bpf, and cgroups v2 boundaries.
\end{enumerate}

\subsection{Formal Distributed Audit Model}

\subsubsection{Audit Invariants}
Audit processing in ERP.1 satisfies four fundamental invariants:
\begin{enumerate}
    \item \textbf{Irreversibility}: Audit entries form an append-only cryptographic hash chain linked via SHA-384 hashes. Past entries cannot be modified or reordered.
    \item \textbf{Non-Repudiation}: Critical records and batch checkpoints are signed by Security Admin BEAM using ECDSA P-384 keys and sealed with RFC 3161 timestamps.
    \item \textbf{Isolation}: Application BEAMs cannot read or mutate audit structures of higher planes.
    \item \textbf{Completeness}: No privileged action (certificate issuance, key unsealing, ABAC policy change, device grant) can execute without a synchronous audit record commit (fail-secure, AU-5).
\end{enumerate}

\subsubsection{Data Models (\texttt{\#audit\_record\{\}} and \texttt{\#audit\_checkpoint\{\}})}
Every audit event follows the mandatory schema specified by NIST AU-3:
\begin{lstlisting}[language=Erlang]
-type audit_event() ::
    auth_login | auth_logout | auth_failure | cert_issue | cert_revoke |
    key_gen | key_rotate | key_unseal | policy_change | privilege_escalation |
    sys_boot | sys_shutdown | process_spawn | process_exit | fs_mount |
    device_grant | net_bind | config_update | user_action | data_access |
    data_mutation | session_start | session_stop | api_call | atom().

-type audit_action() ::
    create | read | update | delete | execute | authenticate | authorize |
    grant | revoke | sign | verify | encrypt | decrypt | unseal | rotate |
    mount | unmount | start | stop | atom().

-record(audit_record, {
    id        :: binary(),             % Monotonic ID (64-bit TS + Rand)
    ts        :: integer(),            % Milliseconds UTC / HLC
    plane     :: security | system | application,
    subject   :: binary(),             % Subject identity
    event     :: audit_event(),        % Formalized event atom
    resource  :: binary(),             % Resource identifier
    action    :: audit_action(),       % Formalized action atom
    outcome   :: success | failure,    % Result
    meta = #{}:: map(),                % Metadata context
    prev_hash :: binary(),             % SHA-384 of previous link
    hmac      :: binary(),             % HMAC-SHA-512 (System plane)
    sig       :: binary() | undefined, % ECDSA P-384 (Security plane)
    tsp       :: binary() | undefined  % RFC 3161 timestamp
}).
\end{lstlisting}

Periodic Merkle checkpoints and log segment seals utilize the \texttt{\#audit\_checkpoint\{\}} record schema:
\begin{lstlisting}[language=Erlang]
-record(audit_checkpoint, {
    node_id      :: binary(),          % Originating node identifier
    from_id      :: binary(),          % First record ID in checkpoint
    to_id        :: binary(),          % Last record ID in checkpoint
    head_hash    :: binary(),          % Top hash of local hash chain
    record_count :: non_neg_integer(), % Total records in segment
    merkle_root  :: binary(),          % SHA-384 Merkle root of records
    sig          :: binary() | undefined, % Security Admin ECDSA P-384 sig
    tsp          :: binary() | undefined  % RFC 3161 timestamp token
}).
\end{lstlisting}

\subsection{Formal Event and Action Taxonomy}

To prevent field ambiguity, support deterministic SIEM log processing, and satisfy NIST AU-2 (Event Selection) and AU-3 (Content of Audit Records), \texttt{synrc/au} formalizes the allowed atom sets for the \texttt{event} and \texttt{action} fields of \texttt{\#audit\_record\{\}}.

\subsubsection{Formalized Audit Events (\texttt{audit\_event()})}
Audit events represent functional domain categories partitioned by execution plane:

\begin{table}[ht!]
\centering
\small
\noindent\makebox[\textwidth][c]{
\begin{tabular}{llp{7cm}}
\toprule
\textbf{Plane} & \textbf{Event Atom} & \textbf{Semantics \& Triggering Context} \\
\midrule
\textbf{Security} & \texttt{auth\_login} & User or service principal authentication attempt \\
 & \texttt{auth\_logout} & Session termination or revocation \\
 & \texttt{auth\_failure} & Authentication rejection or credential failure \\
 & \texttt{cert\_issue} & X.509 / TLS Certificate issuance by Security CA \\
 & \texttt{cert\_revoke} & Certificate revocation or CRL update \\
 & \texttt{key\_gen} & Cryptographic keypair generation (ECDSA / AES) \\
 & \texttt{key\_rotate} & Key lifecycle rotation or re-keying operation \\
 & \texttt{key\_unseal} & Master key unsealing or vault hardware access \\
 & \texttt{policy\_change} & ABAC / RBAC security policy modification \\
 & \texttt{privilege\_escalation} & Elevating capability or ambient authority \\
\midrule
\textbf{System} & \texttt{sys\_boot} & Operating system or kernel initialization \\
 & \texttt{sys\_shutdown} & System shutdown, reboot, or emergency halt \\
 & \texttt{process\_spawn} & BEAM process or container execution spawn \\
 & \texttt{process\_exit} & Process termination or abnormal crash \\
 & \texttt{fs\_mount} & Storage volume / Landlock directory mount \\
 & \texttt{device\_grant} & Hardware seat or device capability assignment \\
 & \texttt{net\_bind} & Network socket bind or outbound connection \\
 & \texttt{config\_update} & Kernel parameter or system config update \\
\midrule
\textbf{Application} & \texttt{user\_action} & User interface or workflow activity \\
 & \texttt{data\_access} & Database read or query operation \\
 & \texttt{data\_mutation} & Record insertion, update, or deletion \\
 & \texttt{session\_start} & Application domain session initialization \\
 & \texttt{session\_stop} & Application domain session destruction \\
 & \texttt{api\_call} & Service API invocation or RPC message \\
\bottomrule
\end{tabular}
}
\caption{Formal Taxonomy of Allowed Audit Event Atoms.}
\end{table}

\subsubsection{Formalized Audit Actions (\texttt{audit\_action()})}
Audit actions specify the discrete operational action executed on a targeted resource:

\begin{table}[ht!]
\centering
\small
\noindent\makebox[\textwidth][c]{
\begin{tabular}{lp{10cm}}
\toprule
\textbf{Action Atom} & \textbf{Operational Semantics} \\
\midrule
\texttt{create} & Instantiation of a new record, entity, or resource \\
\texttt{read} & Accessing, inspecting, or querying resource content \\
\texttt{update} & Modifying or mutating existing resource state \\
\texttt{delete} & Removing, purging, or soft-deleting a resource \\
\texttt{execute} & Launching or invoking an executable, task, or process \\
\texttt{authenticate} & Verifying identity credentials \\
\texttt{authorize} & Evaluating access control capabilities or permissions \\
\texttt{grant} & Bestowing privileges, roles, or capability tokens \\
\texttt{revoke} & Rescinding privileges, roles, or capability tokens \\
\texttt{sign} & Calculating cryptographic digital signatures \\
\texttt{verify} & Validating signatures, Merkle proofs, or hashes \\
\texttt{encrypt} & Encrypting data at rest or in transit \\
\texttt{decrypt} & Decrypting encrypted payloads \\
\texttt{unseal} & Accessing sealed cryptographic keys or vaults \\
\texttt{rotate} & Replacing cryptographic keys or certificates \\
\texttt{mount} & Mounting filesystems or capability trees \\
\texttt{unmount} & Unmounting filesystems or capability trees \\
\texttt{start} & Launching a service, node, or application process \\
\texttt{stop} & Terminating a service, node, or application process \\
\bottomrule
\end{tabular}
}
\caption{Formal Taxonomy of Allowed Audit Action Atoms.}
\end{table}

\subsubsection{Multi-Node CRDT Log Merge}
In multi-node / multi-DC environments, each node maintains a local sequential hash chain. Global consensus is achieved using an \textbf{Add-only Merkle-CRDT} (G-Set of signed checkpoints):
\begin{equation}
\text{GlobalLog} = \text{G-Set}\left( \{ \text{NodeId}, \text{HeadHash}, \text{SignedCheckpoint}, \text{MerkleRoot} \} \right)
\end{equation}
Merging two node logs $L_A$ and $L_B$ is deterministic and conflict-free:
\begin{equation}
\text{merge}(L_A, L_B) = L_A \cup L_B
\end{equation}
Linearization sorts records by the tuple $(TS, \text{NodeId}, \text{RecordId})$, producing a deterministic global sequence without altering original node hash chains.

\newpage
\subsection{NIST Cryptographic Profile}
All cryptographic primitives are standard FIPS/NIST compliant:
\begin{table}[ht!]
\centering
\begin{tabular}{llll}
\toprule
\textbf{Purpose} & \textbf{Algorithm} & \textbf{Specification}  & \textbf{NIST} \\
\midrule
Chain & Merkle Hashing & SHA-384 (or SHA-512) & NIST FIPS 180-4 \\
Message Integrity & HMAC-SHA-512 & NIST FIPS 198-1 \\
Digital Signature & ECDSA P-384 & NIST FIPS 186-4 (secp384r1) \\
Timestamp & RFC 3161 & SHA-384 / SHA-512 TimeStampToken \\
Key Generation & ECDH P-384 & NIST SP 800-56A \\
\bottomrule
\end{tabular}
\caption{NIST Cryptographic Profile for \texttt{synrc/audit}.}
\end{table}

\subsection{NIST SP 800-53 Rev. 5 AU Control Attestation}

\begin{table}[t!]
\centering
\small
\noindent\makebox[\textwidth][c]{
\begin{tabular}{lp{5cm}c p{7cm}}
\toprule
\textbf{Control} & \textbf{Description} & \textbf{Plane} & \textbf{Implementation in \texttt{synrc/audit}} \\
\midrule
AU-1 & Audit Policy and Procedures & [C] & Defined in CMDB metadata. \\
AU-2 & Event Selection & [C],[A] & Configurable event masks per plane. \\
AU-3 & Content of Audit Records & [C] & Mandatory \texttt{\#audit\_record\{\}} schema. \\
AU-4 & Storage Capacity & [C] & ETS / KVS sizing and segmentation. \\
AU-5 & Response to Processing Failures & [C] & Synchronous fail-secure \texttt{gen\_server} call. \\
AU-6 & Review, Analysis, Reporting & [C],[S] & Offline verifier (\texttt{audit\_verify}) \& SIEM export. \\
AU-7 & Audit Reduction & [C] & Filtered queries and Merkle proofs. \\
AU-8 & Time Stamps & [C],[S] & UTC ms / HLC + RFC 3161 timestamps. \\
AU-9 & Protection of Audit Info & [C],[S] & Append-only SHA-384 chain + HMAC-SHA-512. \\
AU-10 & Non-Repudiation & [S] & ECDSA P-384 signatures by Security plane. \\
AU-11 & Audit Record Retention & [C] & Archiving tool (\texttt{audit\_archive}) for cold storage. \\
AU-12 & Audit Generation & [A],[C] & Real-time generation via capability bridge. \\
\bottomrule
\end{tabular}
}
\caption{NIST SP 800-53 Rev. 5 AU Controls Attestation Matrix.}
\end{table}

\subsection{Comprehensive Audit Usage \& Code Examples}

\subsubsection{1. Starting the Audit Service}
To initialize the System Admin BEAM aggregator with ECDSA P-384 Security keys:
\begin{lstlisting}[language=Erlang]
{PubKey, PrivKey} = audit_crypto:generate_keypair(),
MacKey = audit_crypto:generate_mac_key(),

{ok, Pid} = audit_core:start_link(#{
     node_id     => <<"DC1-Node01">>,
     boot_time   => erlang:system_time(millisecond),
     mac_key     => MacKey,
     sec_keypair => {PubKey, PrivKey}}).
\end{lstlisting}

\subsubsection{2. Logging Standard Application Events (AU-12)}
From any Application plane process:
\begin{lstlisting}[language=Erlang]
ok = audit_client:log(
     application,
     <<"user_session_42">>,
     user_login,
     <<"/api/v1/auth">>,
     authenticate,
     success,
     #{<<"ip">> => <<"192.168.1.100">>, <<"seat">> => <<"seat0">>}).
\end{lstlisting}

\newpage
\subsubsection{3. Logging Critical Security Events with ECDSA Signatures}
For high-assurance operations requiring non-repudiation (AU-10):
\begin{lstlisting}[language=Erlang]
ok = audit_client:log_critical(
     security,
     <<"security_admin">>,
     cert_issue,
     <<"CA-Root-2026">>,
     sign,
     success,
     #{<<"subject_dn">> => <<"CN=SystemAdmin">>}).
\end{lstlisting}

\subsubsection{4. Generating and Verifying Merkle Checkpoints}
To generate a signed checkpoint for active records:
\begin{lstlisting}[language=Erlang]
Checkpoint = audit_core:get_checkpoint(),
Valid      = audit_checkpoint:verify(Checkpoint, PubKey),
io:format("Checkpoint valid: ~p~n", [Valid]).
\end{lstlisting}

\subsubsection{5. Multi-Node CRDT Log Merging, Delta \& Sync}
Merging, diffing, and linearizing audit streams from independent nodes:
\begin{lstlisting}[language=Erlang]
SetA          = audit_merge:add_node_log(audit_merge:new_log_set(), CheckpointA, RecordsA),
SetB          = audit_merge:add_node_log(audit_merge:new_log_set(), CheckpointB, RecordsB),
MergedSet     = audit_merge:merge(SetA, SetB),
true          = audit_merge:verify_merged(MergedSet, PubKey),
LinearRecords = audit_merge:linearize(MergedSet, #{plane => security, limit => 50}),

%% Calculate synchronization delta and missing entries diff
DiffSet       = audit_merge:diff(SetA, SetB),
DeltaSet      = audit_merge:delta(MergedSet, #{<<"DC1-Node01">> => 10}),
Stats         = audit_merge:stats(MergedSet).
\end{lstlisting}

\subsubsection{6. Pure Offline Verification (\texttt{audit\_verify})}
Perform independent offline verification of a record chain and Merkle inclusion:
\begin{lstlisting}[language=Erlang]
GenesisHash = audit_chain:genesis_hash(<<"DC1-Node01">>, BootTime),
case audit_verify:verify_chain(Records, GenesisHash, MacKey, PubKey) of
    ok -> io:format("Chain integrity verified.~n");
    {error, {tampered_record, Index, BadRecord}} ->
        io:format("Tampering detected at record ~p!~n", [Index])
end,
IsIncluded = audit_verify:verify_inclusion(Record, Checkpoint).
\end{lstlisting}

\subsubsection{7. Cold Storage Log Archiving \& Sealing (\texttt{audit\_archive})}
Splitting, sealing, and serializing log segments for long-term cold storage (AU-11):
\begin{lstlisting}[language=Erlang]
{Segment, Remaining} = audit_archive:cut_segment(Records, CutoffTimestamp),
SealedCP             = audit_archive:seal_segment(Segment, <<"DC1-Node01">>, PrivKey),
ArchiveBin           = audit_archive:serialize_archive(Segment, SealedCP),

{ok, RestoredCP, RestoredRecords} = audit_archive:deserialize_archive(ArchiveBin).
\end{lstlisting}

\subsubsection{8. SIEM \& OSCAL Export (\texttt{audit\_export})}
Exporting log entries for Wazuh / Syslog or compliance evidence:
\begin{lstlisting}[language=Erlang]
CEFString = audit_export:to_siem(Record),
JSONData  = audit_export:to_json(Records),
OscalMap  = audit_export:to_oscal(Checkpoint, Records).
\end{lstlisting}

\subsection{Erlang Module API Reference}
This section provides the complete reference specification for all exported functions across the 12 modules in \texttt{synrc/au}.

\subsubsection{\texttt{audit\_client}}
High-level Application Plane client interface for emitting audit events to \texttt{audit\_core}.
\begin{itemize}
    \item \texttt{log(Subject, Event, Resource, Action, Outcome, Meta) -> ok | \{error, term()\}} \\
    Submits a standard audit event defaulting to the \texttt{application} plane.
    \item \texttt{log(Plane, Subject, Event, Resource, Action, Outcome, Meta) -> ok | \{error, term()\}} \\
    Submits a standard audit event targeting the specified execution plane (\texttt{security}, \texttt{system}, or \texttt{application}).
    \item \texttt{log\_critical(Plane, Subject, Event, Resource, Action, Outcome, Meta) -> ok | \{error, term()\}} \\
    Submits a high-assurance audit event requiring Security Admin ECDSA P-384 signature and RFC 3161 timestamping.
\end{itemize}

\subsubsection{\texttt{audit\_core}}
OTP \texttt{gen\_server} managing local ETS storage, append-only hash chain linking, and checkpoint generation.
\begin{itemize}
    \item \texttt{start\_link() -> \{ok, pid()\} | \{error, term()\}} \\
    Starts the audit core server with default node options.
    \item \texttt{start\_link(Opts::map()) -> \{ok, pid()\} | \{error, term()\}} \\
    Starts the audit core server with custom options (\texttt{node\_id}, \texttt{boot\_time}, \texttt{mac\_key}, \texttt{sec\_keypair}).
    \item \texttt{stop() -> ok} \\
    Gracefully stops the \texttt{audit\_core} server.
    \item \texttt{log\_event(Record::\#audit\_record\{\}) -> ok | \{error, term()\}} \\
    Synchronously appends a standard audit record to the local hash chain and ETS table.
    \item \texttt{log\_critical\_event(Record::\#audit\_record\{\}) -> ok | \{error, term()\}} \\
    Synchronously appends a signed critical audit record to the local hash chain and ETS table.
    \item \texttt{get\_head\_hash() -> binary()} \\
    Returns the current 384-bit SHA-384 head hash of the node's local audit chain.
    \item \texttt{get\_records() -> [\#audit\_record\{\}]} \\
    Retrieves all audit records stored in the local ETS table in chronological sequence.
    \item \texttt{get\_checkpoint() -> \#audit\_checkpoint\{\}} \\
    Generates a signed \texttt{\#audit\_checkpoint\{\}} spanning all currently accumulated records.
    \item \texttt{set\_security\_keys(PubKey::binary(), PrivKey::binary()) -> ok} \\
    Dynamically configures or updates the Security Admin ECDSA P-384 keypair.
\end{itemize}

\subsubsection{\texttt{audit\_record}}
Constructs, validates, encodes, and formats individual NIST AU-3 audit record instances.
\begin{itemize}
    \item \texttt{new(Plane, Subject, Event, Resource, Action, Outcome, Meta) -> \#audit\_record\{\}} \\
    Instantiates and validates a new \texttt{\#audit\_record\{\}} with automatic monotonic ID generation.
    \item \texttt{allowed\_events() -> [audit\_event()]} \\
    Returns the list of standard formalized audit event atoms.
    \item \texttt{allowed\_actions() -> [audit\_action()]} \\
    Returns the list of standard formalized audit action atoms.
    \item \texttt{is\_allowed\_event(Event::atom()) -> boolean()} \\
    Checks whether a given atom belongs to the formalized \texttt{audit\_event()} set.
    \item \texttt{is\_allowed\_action(Action::atom()) -> boolean()} \\
    Checks whether a given atom belongs to the formalized \texttt{audit\_action()} set.
    \item \texttt{validate(Record::term()) -> ok | \{error, term()\}} \\
    Enforces NIST AU-3 schema validation on record fields.
    \item \texttt{canonical\_payload(\#audit\_record\{\}) -> binary()} \\
    Computes a deterministic binary payload encoding (excluding transient hashes/signatures) used for hashing and MAC generation.
    \item \texttt{encode(\#audit\_record\{\}) -> binary()} \\
    Encodes an audit record into a deterministic binary term format.
    \item \texttt{decode(Bin::binary()) -> \{ok, \#audit\_record\{\}\} | \{error, term()\}} \\
    Decodes a binary term into a validated \texttt{\#audit\_record\{\}}.
    \item \texttt{new\_id() -> binary()} \\
    Generates a monotonic 16-byte record identifier (64-bit UTC ms timestamp + 8-byte random entropy).
    \item \texttt{now\_ms() -> integer()} \\
    Returns current UTC timestamp in milliseconds.
\end{itemize}

\subsubsection{\texttt{audit\_chain}}
Implements append-only hash chain linking, HMAC tag calculation, and link verification.
\begin{itemize}
    \item \texttt{genesis\_hash(NodeId::binary(), BootTime::integer()) -> binary()} \\
    Computes the root genesis hash for a node audit chain (\texttt{ERP.1-AUDIT-ROOT:NodeId:BootTime}).
    \item \texttt{append\_record(Record, PrevHash, MacKey) -> \#audit\_record\{\}} \\
    Calculates \texttt{prev\_hash} via SHA-384 and \texttt{hmac} via HMAC-SHA-512 over \texttt{(PrevHash || Payload)}.
    \item \texttt{append\_critical\_record(Record, PrevHash, MacKey, SecPrivKey) -> \#audit\_record\{\}} \\
    Appends standard chain hashes and attaches an ECDSA P-384 signature and RFC 3161 timestamp token.
    \item \texttt{verify\_record(Record, PrevPrevHash, MacKey, SecPubKey) -> boolean()} \\
    Validates single record hash continuity, HMAC tag, ECDSA signature, and timestamp token.
\end{itemize}

\subsubsection{\texttt{audit\_checkpoint}}
Constructs and verifies Merkle checkpoints over log segments.
\begin{itemize}
    \item \texttt{create(NodeId, Records, HeadHash, SecPrivKey) -> \#audit\_checkpoint\{\}} \\
    Builds a \texttt{\#audit\_checkpoint\{\}}, computes the Merkle tree root over record canonical payloads, and applies an optional ECDSA signature.
    \item \texttt{verify(Checkpoint::\#audit\_checkpoint\{\}, SecPubKey) -> boolean()} \\
    Validates ECDSA P-384 signature and RFC 3161 timestamp token on a checkpoint.
    \item \texttt{compute\_merkle\_root([\#audit\_record\{\}]) -> binary()} \\
    Computes the binary Merkle root hash (SHA-384) over a list of audit records.
\end{itemize}

\subsubsection{\texttt{audit\_merge}}
Add-only Merkle-CRDT G-Set engine for multi-node log synchronization, diffing, delta tracking, and linearization.
\begin{itemize}
    \item \texttt{new\_log\_set() -> log\_set()} \\
    Returns an empty G-Set map.
    \item \texttt{add\_node\_log(LogSet, Checkpoint, Records) -> log\_set()} \\
    Adds or updates a node log entry in the G-Set, merging records if the node exists.
    \item \texttt{merge(LogSetA, LogSetB) -> log\_set()} \\
    Computes the deterministic CRDT G-Set union of two log sets.
    \item \texttt{linearize(LogSet) -> [\#audit\_record\{\}]} \\
    Produces a total order sequence of all records across nodes sorted by \texttt{(TS, NodeId, RecordId)}.
    \item \texttt{linearize(LogSet, Opts::map()) -> [\#audit\_record\{\}]} \\
    Linearizes records with filtering options (\texttt{from\_ts}, \texttt{to\_ts}, \texttt{plane}, \texttt{limit}).
    \item \texttt{global\_merkle\_root(LogSet) -> binary()} \\
    Computes a cluster-wide global Merkle root over sorted node head hashes.
    \item \texttt{verify\_merged(LogSet, SecPubKey) -> boolean()} \\
    Verifies all checkpoints and record hash chains contained within a merged log set.
    \item \texttt{diff(LogSetA, LogSetB) -> log\_set()} \\
    Calculates missing records present in \texttt{LogSetB} but absent in \texttt{LogSetA}.
    \item \texttt{delta(LogSet, KnownCounts::map()) -> log\_set()} \\
    Generates a delta log set containing only records newer than specified per-node counts.
    \item \texttt{apply\_delta(LocalSet, DeltaSet) -> log\_set()} \\
    Applies a delta log set into a local log set via CRDT merge.
    \item \texttt{stats(LogSet) -> map()} \\
    Computes aggregate cluster metrics (\texttt{node\_count}, \texttt{total\_records}, \texttt{min\_timestamp}, \texttt{max\_timestamp}, \texttt{global\_merkle\_root}).
\end{itemize}

\subsubsection{\texttt{audit\_verify}}
Pure offline cryptographic chain and checkpoint verifier.
\begin{itemize}
    \item \texttt{verify\_chain(Records, GenesisHash, MacKey, SecPubKey) -> ok | \{error, \{tampered\_record, integer(), term()\}\}} \\
    Iteratively verifies chain continuity, HMAC tags, and ECDSA signatures from genesis.
    \item \texttt{verify\_checkpoint(Checkpoint, SecPubKey) -> boolean()} \\
    Verifies a checkpoint's signature and timestamp offline.
    \item \texttt{verify\_inclusion(Record, Checkpoint) -> boolean()} \\
    Verifies whether a record payload hash is included in a checkpoint's Merkle root.
\end{itemize}

\subsubsection{\texttt{audit\_archive}}
Cold storage segmentation, sealing, and serialization primitives.
\begin{itemize}
    \item \texttt{cut\_segment(Records, CutoffTS) -> \{Segment, Remaining\}} \\
    Partitions audit records at a timestamp boundary.
    \item \texttt{seal\_segment(Segment, NodeId, SecPrivKey) -> \#audit\_checkpoint\{\}} \\
    Seals a record segment into a signed \texttt{\#audit\_checkpoint\{\}}.
    \item \texttt{serialize\_archive(Segment, Checkpoint) -> binary()} \\
    Compresses and serializes a segment and its checkpoint into binary format.
    \item \texttt{deserialize\_archive(Bin::binary()) -> \{ok, \#audit\_checkpoint\{\}, [\#audit\_record\{\}]\} | \{error, term()\}} \\
    Deserializes and validates an archive binary blob.
\end{itemize}

\subsubsection{\texttt{audit\_export}}
Formats audit evidence into standard SIEM, JSON, and OSCAL representations.
\begin{itemize}
    \item \texttt{to\_json(RecordOrRecords) -> binary()} \\
    Formats record(s) into structured JSON.
    \item \texttt{to\_siem(Record) -> binary()} \\
    Formats an audit record into Wazuh / Syslog Common Event Format (CEF).
    \item \texttt{to\_oscal(Checkpoint, Records) -> map()} \\
    Formats checkpoint and records into a NIST OSCAL Assessment Results structure.
\end{itemize}

\subsubsection{\texttt{audit\_crypto}}
FIPS/NIST cryptographic primitives (ECDSA P-384, SHA-384, HMAC-SHA-512, RFC 3161 timestamps).
\begin{itemize}
    \item \texttt{generate\_keypair() -> \{PubKey, PrivKey\}} \\
    Generates a fresh ECDSA P-384 (\texttt{secp384r1}) keypair.
    \item \texttt{generate\_mac\_key() -> binary()} \\
    Generates a 256-bit symmetric key for HMAC calculation.
    \item \texttt{hash(Data::binary()) -> binary()} \\
    Computes SHA-384 hash of binary data.
    \item \texttt{hash(Algo, Data::binary()) -> binary()} \\
    Computes specified digest (\texttt{sha384} or \texttt{sha512}).
    \item \texttt{hmac(Key, Data) -> binary()} \\
    Computes HMAC-SHA-512 over data using symmetric key.
    \item \texttt{sign(PrivKey, Data) -> binary()} \\
    Signs data digest using ECDSA P-384 with SHA-384 digest.
    \item \texttt{verify(PubKey, Data, Signature) -> boolean()} \\
    Verifies ECDSA P-384 signature over data.
    \item \texttt{timestamp\_token(Digest, TS) -> binary()} \\
    Generates a simulated RFC 3161 timestamp token for a given digest and timestamp.
    \item \texttt{verify\_timestamp(Token, Digest, TS) -> boolean()} \\
    Verifies an RFC 3161 timestamp token.
\end{itemize}

\subsubsection{\texttt{audit\_app} \& \texttt{audit\_sup}}
OTP application lifecycle and supervisor specification.
\begin{itemize}
    \item \texttt{audit\_app:start(StartType, StartArgs) -> \{ok, pid()\} | \{error, term()\}} \\
    Starts the \texttt{au} OTP application and supervisor tree.
    \item \texttt{audit\_app:stop(State) -> ok} \\
    Stops the application.
    \item \texttt{audit\_sup:start\_link() -> \{ok, pid()\} | \{error, term()\}} \\
    Starts the \texttt{audit\_sup} supervisor (\texttt{one\_for\_one} strategy supervising \texttt{audit\_core}).
    \item \texttt{audit\_sup:init(Args) -> \{ok, \{SupFlags, [ChildSpec]\}\}} \\
    OTP supervisor callback specifying child worker parameters.
\end{itemize}

\section{Related Work, Historical Lineage, and Ecosystem Comparison}

\subsection{Historical Lineage}
Cryptographic audit logging and tamper-evident event recording rest upon foundational work across distributed security literature:
\begin{enumerate}
    \item \textbf{Schneier \& Kelsey Forward-Secure Logging (1998, 1999)}~\cite{schneier1998,schneier1999}: Schneier and Kelsey established the core model for secure audit logs on untrusted machines using evolving hash chains ($H_i = \text{SHA}(H_{i-1} \parallel W_i)$) and forward-secure symmetric MAC keys deleted after writing. \texttt{synrc/au} builds upon this linear chain model, extending it to a multi-plane Erlang execution architecture where key generation and signature authorization reside strictly inside the isolated Security Admin BEAM [S].
    \item \textbf{Merkle Trees \& Audit Proofs (Merkle 1987, Crosby \& Wallach 2009)}~\cite{merkle1987,crosby2009}: Merkle's 1987 binary tree structures and Crosby \& Wallach's efficient data structures for tamper-evident logging established logarithmic $\mathcal{O}(\log N)$ inclusion verification and consistency proofs. \texttt{synrc/au} utilizes SHA-384 Merkle trees inside \texttt{audit\_checkpoint} and for multi-DC global cluster root generation (\texttt{global\_merkle\_root/1}).
    \item \textbf{Certificate Transparency \& RFC 6962 (Laurie et al., 2013)}~\cite{rfc6962}: Established publicly auditable append-only Merkle tree logs for public TLS certificate issuance. \texttt{synrc/au} incorporates append-only Merkle tree semantics while providing native AEAD AES-256-GCM encryption at rest (SC-28) for enterprise confidentiality.
    \item \textbf{Google Transparency Dev Initiative}~\cite{transparencydev}: Modern zero-dependency Go ecosystem (\url{https://transparency.dev}) provides modular, tile-based Merkle tree storage (\texttt{tessera}) and witness verification networks for application transparency. \texttt{synrc/au} shares this zero-dependency philosophy, implementing native Erlang/OTP cryptographic log storage and verification.
    \item \textbf{NIST SP 800-53 Rev. 5 Standards (2020)}~\cite{nist80053}: Outlines mandatory security controls for US Federal and FedRAMP High baseline environments.
\end{enumerate}

\subsection{Modern Zero-Dependency Competitor Analysis Across Languages}

We analyze \texttt{synrc/au} against modern, standalone, minimal-dependency audit and transparency logging implementations in Go, C99, and Rust:

\begin{table}[ht!]
\centering
\scriptsize
\noindent\makebox[\textwidth][c]{
\begin{tabular}{lp{1.8cm}p{1.8cm}p{1.8cm}p{2.2cm}p{4.5cm}}
\toprule
\textbf{System} & \textbf{Language / TCB} & \textbf{Dependency Model} & \textbf{Isolation Architecture} & \textbf{Cryptographic Security} & \textbf{Replication / Merging} \\
\midrule
\textbf{\texttt{synrc/au}} & Erlang/OTP & OTP \texttt{crypto} & Process boundary ([S],[C],[A]) & SHA-384 Chain + ECDSA P-384 + AES-256-GCM AEAD & \textbf{Add-only Merkle-CRDT G-Set} \\
\textbf{Linux \texttt{kauditd}} & C99 Kernel & POSIX Kernel & Ring 0 Kernel vs Ring 3 Userspace & Plaintext sequential logs (Requires external syslog) & Manual syslog forwarder \\
\textbf{\texttt{transparency.dev}} & Go (Tessera) & Go stdlib & Process boundary & Tile Merkle Tree + Witness Signatures & Static HTTP tiles / Checkpoint witnesses \\
\textbf{\texttt{CT-Go}} & Go & Go stdlib & Process boundary & RFC 6962 Merkle Tree + Ed25519 & Centralized log server \\
\textbf{\texttt{Rekor}} & Go & Trillian, gRPC & Container / Microservice & Merkle Transparency Tree + Rekor Sig & DB-backed consensus \\
\textbf{\texttt{tpm2-tss-audit}} & C99, \texttt{libcrypto}, TPM2 & Hardware TPM2 Enclave & TPM2 HMAC / PCR event log signatures & Single node hardware log \\
\textbf{\texttt{sunlight}} & Rust \texttt{std}, \texttt{ring} & Memory-safe process & RFC 6962 Merkle Tree + Ed25519 & Static tile log storage \\
\bottomrule
\end{tabular}
}
\caption{Cross-Language Zero-Dependency Audit System Comparison.}
\end{table}

\subsection{Key Architectural Differentiators of \texttt{synrc/au}}
\begin{enumerate}
    \item \textbf{Three-Plane Process Isolation}: C99 system loggers (\texttt{auditd}) rely on monolithic POSIX processes, exposing signing keys to local root compromise. \texttt{synrc/au} isolates key generation and signing inside the Security Admin BEAM [S], preventing Application [A] or Control [C] planes from accessing long-term private keys.
    \item \textbf{Confidentiality at Rest (SC-28 / AU-9(3))}: Standard transparency loggers (Rekor, CT-Go, \texttt{transparency.dev}) are public plaintext ledgers. \texttt{synrc/au} natively encrypts payload metadata using AES-256-GCM AEAD authenticated encryption.
    \item \textbf{Asynchronous Conflict-Free Merging}: Uses Add-only Merkle-CRDT G-Set semantics to merge logs across multi-DC deployments deterministically without centralized database locks.
\end{enumerate}

\section{Conclusion}

I dedicate this work to my daugther Michelle.

The \texttt{synrc/audit} library provides a minimal, self-contained, and cryptographically
verifiable audit infrastructure for ERP.1 systems. By combining standard NIST algorithms,
Erlang/OTP supervision, three-plane capability separation, CRDT Merkle log merging,
cold storage log archiving, and comprehensive export capabilities, it fulfills the full
spectrum of NIST SP 800-53 Rev. 5 AU requirements without external dependencies.

\newpage
\begin{thebibliography}{99}

\bibitem{schneier1998}
B.~Schneier and J.~Kelsey,
``Cryptographic support for secure logs on untrusted machines,''
\emph{ACM Transactions on Information and System Security (TISSEC)}, vol.~2, no.~2, pp.~159--176, 1998.

\bibitem{schneier1999}
B.~Schneier and J.~Kelsey,
``Secure audit logs supporting tamper detection and forward-secrecy,''
in \emph{Proceedings of the 10th USENIX Security Symposium}, 1999, pp.~139--165.

\bibitem{merkle1987}
R.~C. Merkle,
``A digital signature based on a conventional encryption function,''
in \emph{Conference on the Theory and Application of Cryptographic Techniques (CRYPTO '87)}, Springer, 1987, pp.~369--378.

\bibitem{crosby2009}
S.~A. Crosby and D.~S. Wallach,
``Efficient data structures for tamper-evident logging,''
in \emph{Proceedings of the 18th USENIX Security Symposium}, 2009, pp.~317--334.

\bibitem{rfc6962}
B.~Laurie, A.~Langley, and E.~Kasper,
``Certificate Transparency,''
\emph{IETF RFC 6962}, Jun. 2013.

\bibitem{transparencydev}
Google Transparency Engineering,
``Verifiable Data Structures and Transparency Logs,''
\emph{Google Transparency Dev Initiative}, 2024. [Online]. Available: \url{https://transparency.dev}

\bibitem{rfc3161}
C.~Adams, P.~Cain, D.~Pinkas, and R.~Zuccherato,
``Internet X.509 Public Key Infrastructure Time-Stamp Protocol (TSP),''
\emph{IETF RFC 3161}, Aug. 2001.

\bibitem{nist80053}
National Institute of Standards and Technology (NIST),
``Security and Privacy Controls for Information Systems and Organizations,''
\emph{NIST Special Publication 800-53}, Revision 5, Sep. 2020.

\bibitem{fips186}
National Institute of Standards and Technology (NIST),
``Digital Signature Standard (DSS),''
\emph{FIPS PUB 186-4}, Jul. 2013.

\bibitem{fips180}
National Institute of Standards and Technology (NIST),
``Secure Hash Standard (SHS),''
\emph{FIPS PUB 180-4}, Aug. 2015.

\bibitem{fips198}
National Institute of Standards and Technology (NIST),
``The Keyed-Hash Message Authentication Code (HMAC),''
\emph{FIPS PUB 198-1}, Jul. 2008.

\bibitem{belllapadula}
D.~E. Bell and L.~J. LaPadula,
``Secure Computer System: Unified Exposition and Multics Interpretation,''
\emph{MITRE Corp Technical Report MTR-2997}, Jul. 1975.

\bibitem{shapiro2011}
M.~Shapiro, N.~Pregui\c{c}a, C.~Baquero, and M.~Zawirski,
``Conflict-free Replicated Data Types,''
in \emph{Symposium on Self-Stabilizing Systems (SSS 2011)}, Springer, 2011, pp.~386--400.

\bibitem{kauditd}
Linux Audit Subsystem Project,
``Linux Audit Framework (\texttt{kauditd} \& \texttt{auditd}),''
\emph{Linux Kernel Documentation}, 2024. [Online]. Available: \url{https://docs.kernel.org/trace/events.html}

\bibitem{rekor}
Sigstore Community,
``Rekor: Signature Transparency Log for Software Supply Chain Security,''
\emph{Linux Foundation Open Source Project}, 2022. [Online]. Available: \url{https://www.sigstore.dev}

\bibitem{trillian}
Google Transparency Team,
``Trillian: Verifiable Data Structures for Verifiable Logs,''
\emph{GitHub Repository}, 2023. [Online]. Available: \url{https://github.com/google/trillian}

\end{thebibliography}

\end{document}
